%!LW recipe=xelatex
\documentclass[aspectratio=169]{beamer}

\usepackage{tikz}

\usepackage{amsmath}
\usepackage{tikz-network}
\usepackage{tikz-cd}
\usepackage{tikzit}
\usepackage{subcaption}
\usepackage{minted}
\usepackage{mathtools}
\usepackage{stmaryrd}
\newcommand{\consistency}{{\smile}}
\usepackage{adjustbox}
\usepackage{appendixnumberbeamer}
\usepackage{csquotes}
\usepackage{mathpartir}
\usepackage{macros}
\usepackage{longtable}
\usepackage{booktabs}
% for tables
\usepackage{multirow}
\usepackage{bigdelim}
\usepackage{colortbl}
\usepackage{colortbl}
\makeatletter
\providecommand\CT@arc@{}
\providecommand\CT@row@color{}
\providecommand\CT@cell@color{}
\providecommand\CT@do@color{}
\makeatother
\usepackage[style=authoryear,backend=biber]{biblatex} % or backend=biber
\addbibresource{./bibliography.bib}
\setbeamerfont{author}{series=\normalfont}
% \hypersetup{
%     colorlinks=true,
%     linkcolor=gray,
%     filecolor=magenta,      
%     urlcolor=cyan,
%     pdftitle={Overleaf Example},
%     pdfpagemode=FullScreen,
%     }

\input{../../../sample.tikzstyles}
\input{../../../hypergraph.tikzstyles}
\input{../../../hypergraph.tikzdefs}

\usetheme[subsectionpage=simple]{metropolis}
% \usetheme{awesome}

% \setsansfont{Ubuntu}
% \setmonofont{Ubuntu Mono}


\setbeamertemplate{footline}
{
  \leavevmode
  \hbox{
  \begin{beamercolorbox}[wd=.15\paperwidth,ht=2.25ex,dp=1ex,center]{title in head/foot}
    \usebeamerfont{author in head/foot}\insertshortauthor
  \end{beamercolorbox}

  \begin{beamercolorbox}[wd=.7\paperwidth,ht=2.25ex,dp=1ex,center]{author in head/foot}
    \usebeamerfont{author in head/foot}\insertshorttitle
  \end{beamercolorbox}

  \begin{beamercolorbox}[wd=.15\paperwidth,ht=2.25ex,dp=1ex,center]{title in head/foot}
    \insertframenumber{} / \inserttotalframenumber
  \end{beamercolorbox}
  }
}

\makeatletter
\metropolis@disablesubsectionpage
\makeatother

\newcommand{\bsubsection}[1]{\subsection{$\bullet$ #1}}

\title{Categorical E-Graphs for $\lambda$-calculi} 
\date{July, 2025} % Use metropolis theme 
\author[Aleksei Tiurin]{\underline{Aleksei Tiurin}, Dan Ghica, Nick Hu\\University of Birmingham}
\begin{document} 


\renewcommand{\maketitle}{
  \begin{frame}[plain]
    \usebeamerfont{title}\usebeamercolor[fg]{title}\inserttitle\par
    \alert{\rule{\textwidth}{1pt}}
    \vfill
    \[
    \left\llbracket 
        \;\vcenter{\hbox{\tikzfig{./figures/e-graph-logo}}}\; 
    \right\rrbracket
    \]
    \vfill
    \centering
    \usebeamerfont{author}\insertauthor\par
    \usebeamerfont{date}\insertdate\par
  \end{frame}
}

\maketitle 

% \begin{frame}{Table of contents}
%     \small
%     \setbeamertemplate{section in toc}[sections numbered]
%     \tableofcontents
% \end{frame}

\section{Equality saturation $\simeq$ E-graphs}
\bsubsection{Equality Saturation and E-graphs}

\begin{frame}{Equality saturation}
\vfill
First introduced as a decision procedure for the quantifier-free fragment of theory of uninterpreted functional symbols with equality
\pause
\vfill
Program optimisation technique that obviates the need to worry about \alert{phase ordering}
\pause
\vfill
Equalities are applied non-destructively until the representation is \enquote{saturated} with respect to these equalities\par
\pause
\vfill
Once saturated, the \alert{best} candidate can be extracted from the set\par
\vfill
\pause
The \alert{best} candidate is domain-specific
\begin{itemize}
    \item Best program in compiler optimisation
    \item Smallest test-case in fuzzy testing
    \item Smallest digital circuit
    \item $\ldots$
\end{itemize}
\end{frame}

\begin{frame}{E-graphs~[\cite{EggPaper}]}
E-graph is a \alert{data structure} that facilitates equality saturation
\hfill
\begin{example}
\begin{minipage}{0.45\textwidth}
    \[
    \tikzfig{./figures/e-graph-definition-example-2}
    \]
\end{minipage}
\hfill
\begin{minipage}{0.45\textwidth}
E-classes $\scalebox{0.4}{\begin{tikzpicture}\begin{pgfonlayer}{nodelayer}\node [style=empty diag black] (0) at (0, 0) {};\end{pgfonlayer}\end{tikzpicture}}$ 
\begin{itemize}
    \item Set of equivalent terms represented by enodes rooted at this particular e-class
\end{itemize}
\end{minipage}
\end{example}
\end{frame}


\begin{frame}{E-graphs~[\cite{EggPaper}]}
E-graph is a \alert{data structure} that facilitates equality saturation
\hfill
\begin{example}
\begin{minipage}{0.45\textwidth}
    \[
    \tikzfig{./figures/e-graph-definition-example}
    \]
\end{minipage}
\hfill
\begin{minipage}{0.45\textwidth}
E-nodes $\begin{tikzpicture}\begin{pgfonlayer}{nodelayer}\node [style=round box] (0) at (0, 0) {$*$};\node [style=none] (1) at (-0.5, -0.5) {};\node [style=none] (2) at (0.5, -0.5) {};\node [style=none] (3) at (0, -0.5) {$\small\ldots$};\end{pgfonlayer}\begin{pgfonlayer}{edgelayer}\draw [style=new edge style 1] (0.center) to (1.center);\draw [style=new edge style 1] (0.center) to (2.center);\end{pgfonlayer}\end{tikzpicture}$, $\ldots$, $\begin{tikzpicture}\begin{pgfonlayer}{nodelayer}\node [style=round box] (0) at (0, 0) {$<\!\!<$};\node [style=none] (1) at (-0.5, -0.5) {};\node [style=none] (2) at (0.5, -0.5) {};\node [style=none] (3) at (0, -0.5) {$\small\ldots$};\end{pgfonlayer}\begin{pgfonlayer}{edgelayer}\draw [style=new edge style 1] (0.center) to (1.center);\draw [style=new edge style 1] (0.center) to (2.center);\end{pgfonlayer}\end{tikzpicture}$
\begin{itemize}
    \item A representative of an e-class once the root node is fixed
\end{itemize}
\end{minipage}
\end{example}
\end{frame}


\begin{frame}{E-graphs~[\cite{EggPaper}]}
Can also be understood as a device computing the quotient of a free term-algebra over an \alert{algebraic} signature
\pause
    \vfill
    \begin{itemize}
        \item $\Sigma = \{\mathbb{N} \setminus \{0\}, * : 2 \to 1, / : 2 \to 1, <\!\!< : 2 \to 1, a : 0 \to 1\}$
        \item $\mathcal{E} = \{x * 2 = x <\!\!<1, (x * y) / z =  x * (y / z), x / x = 1, x * 1 = x\}$
        \item $a$, $a * 2$, $a <\!\!< 1$, $\ldots$
    \end{itemize}
    \vfill
\begin{example}[$(a * 2) / 2$]
    \[
    \scalebox{0.9}{
    \tikzfig{./figures/e-graph-example-a-no-label}}
    \]
\end{example}
\end{frame}

\begin{frame}{E-graphs~[\cite{EggPaper}]}
    \begin{example}[$x * 2 \to x <\!\!< 1$]
        \hspace{-3em}
        \begin{minipage}{0.25\linewidth}
        \[
        \scalebox{0.8}{
        \tikzfig{./figures/e-graph-example-a-no-label}
        }
        \]    
        \end{minipage}
        \pause
        \hspace{-1.5em}
        \begin{minipage}{0.075\linewidth}
        \[
        \xRightarrow{\texttt{add}(a <\!\!< 1)}
        \]
        \end{minipage}
        \hfill
        \begin{minipage}{0.25\linewidth}
            \[
            \scalebox{0.8}{
            \tikzfig{./figures/e-graph-example-b-add}
            }
            \]    
        \end{minipage}
        \hfill
        \pause
        \begin{minipage}{0.1\linewidth}
            \[
            \xRightarrow{\texttt{merge}(\scalebox{0.3}{\begin{tikzpicture}\begin{pgfonlayer}{nodelayer}\node [style=empty diag yellow] (0) at (0, 0) {<\!\!<};\end{pgfonlayer}\end{tikzpicture}},\scalebox{0.3}{\begin{tikzpicture}\begin{pgfonlayer}{nodelayer}\node [style=empty diag black] (0) at (0, 0) {*};\end{pgfonlayer}\end{tikzpicture}})}
            \]
        \end{minipage}
        \hfill
        \begin{minipage}{0.25\linewidth}
            \[
            \scalebox{0.8}{
            \tikzfig{./figures/e-graph-example-b-merged}
            }
            \]
        \end{minipage}
        \hfill
    \end{example}
\end{frame}


\begin{frame}{E-graphs~[\cite{EggPaper}]}
    \begin{example}[$(x * y) / z \to x * (y / z)$]
        \hspace{-2.5em}
        \begin{minipage}{0.25\linewidth}
        \[
        \scalebox{0.8}{
        \tikzfig{./figures/e-graph-example-b-merged}
        }
        \]    
        \end{minipage}
        \pause
        \hspace{-1.5em}
        \begin{minipage}{0.075\linewidth}
        \[
        \xRightarrow{\texttt{add}(a * (2/2))}
        \]
        \end{minipage}
        \hfill
        \begin{minipage}{0.25\linewidth}
            \[
            \scalebox{0.8}{
                \tikzfig{./figures/e-graph-example-c-add}
            }
            \]    
        \end{minipage}
        \hfill
        \pause
        \begin{minipage}{0.1\linewidth}
            \[
            \xRightarrow{\texttt{merge}(\scalebox{0.3}{\begin{tikzpicture}\begin{pgfonlayer}{nodelayer}\node [style=empty diag yellow] (0) at (0, 0) {/};\end{pgfonlayer}\end{tikzpicture}},\scalebox{0.3}{\begin{tikzpicture}\begin{pgfonlayer}{nodelayer}\node [style=empty diag black] (0) at (0, 0) {*};\end{pgfonlayer}\end{tikzpicture}})}
            \]
        \end{minipage}
        \hfill
        \begin{minipage}{0.25\linewidth}
            \[
            \scalebox{0.8}{
            \tikzfig{./figures/e-graph-example-c-merged}
            }
            \]
        \end{minipage}
        \hfill
    \end{example}
\end{frame}


\begin{frame}{E-graphs~[\cite{EggPaper}]}
    \begin{example}[$(x / x) \to 1$]
        \hspace{-2.5em}
        \begin{minipage}{0.25\linewidth}
        \[
        \scalebox{0.8}{
        \tikzfig{./figures/e-graph-example-c-merged}
        }
        \]    
        \end{minipage}
        \pause
        \hfill
        \begin{minipage}{0.075\linewidth}
        \[
        \xRightarrow{\texttt{add}(1)}
        \]
        \end{minipage}
        \hfill
        \begin{minipage}{0.25\linewidth}
            \[
            \scalebox{0.8}{
            \tikzfig{./figures/e-graph-example-c-merged}
            }
            \]    
        \end{minipage}
        \hfill
        \pause
        \begin{minipage}{0.1\linewidth}
            \[
            \xRightarrow{\texttt{merge}(\scalebox{0.3}{\begin{tikzpicture}\begin{pgfonlayer}{nodelayer}\node [style=empty diag black] (0) at (0, 0) {/};\end{pgfonlayer}\end{tikzpicture}},\scalebox{0.3}{\begin{tikzpicture}\begin{pgfonlayer}{nodelayer}\node [style=empty diag black] (0) at (0, 0) {1};\end{pgfonlayer}\end{tikzpicture}})}
            \]
        \end{minipage}
        \hfill
        \begin{minipage}{0.25\linewidth}
            \[
            \scalebox{0.8}{
            \tikzfig{./figures/e-graph-example-d-merged}
            }
            \]
        \end{minipage}
        \hfill
    \end{example}
\end{frame}


\begin{frame}{E-graphs~[\cite{EggPaper}]}
    \begin{example}[$x * 1 \to x$]
        \hspace{-2em}
        % \hfill
        \begin{minipage}{0.2\linewidth}
        \[
        \scalebox{0.8}{
        \tikzfig{./figures/e-graph-example-d-merged}
        }
        \]    
        \end{minipage}
        \pause
        \hspace{2em}
        \begin{minipage}{0.1\linewidth}
        \[
        \xRightarrow{\texttt{add}(a)}
        \]
        \end{minipage}
        \begin{minipage}{0.2\linewidth}
            \[
            \scalebox{0.8}{
            \tikzfig{./figures/e-graph-example-d-merged}
            }
            \]    
        \end{minipage}
        \hspace{1em}
        \pause
        \begin{minipage}{0.25\linewidth}
            \[
            \xRightarrow{\texttt{merge}(\scalebox{0.3}{\begin{tikzpicture}\begin{pgfonlayer}{nodelayer}\node [style=empty diag black] (0) at (0, 0) {/\;,*};\end{pgfonlayer}\end{tikzpicture}},\scalebox{0.3}{\begin{tikzpicture}\begin{pgfonlayer}{nodelayer}\node [style=empty diag black] (0) at (0, 0) {a};\end{pgfonlayer}\end{tikzpicture}})}
            \]
        \end{minipage}
        \hspace{-2em}
        \begin{minipage}{0.2\linewidth}
            \[
            \scalebox{0.8}{
            \tikzfig{./figures/e-graph-example-e-merged}
            }
            \]
        \end{minipage}
    \end{example}
\end{frame}

\begin{frame}{Challenges}
    \vfill
    \begin{itemize}
    \item Do not scale well in the case of \alert{(symmetric) monoidal theories}
    \item No first-class notion of variables, binding and substitution
    \end{itemize}
    \vfill
\end{frame}

\begin{frame}{Symmetric monoidal theories}
    Presented by a \alert{monoidal signature} $\Sigma = (\Sigma_0, \Sigma_1)$ and a set of \alert{equations} $\mathcal{E}$ between $\Sigma$-terms
    \vfill
    \begin{example}
    \vfill
    $\Sigma = \{ \mu : 1 \to 2, \eta : 1 \to 0\}$
    \vfill
    $\mathcal{E} = \{\mu;\sym \equiv \mu, \mu;(\eta \otimes \mathrm{id}) \equiv \mathrm{id}, \mu;(\mathrm{id} \otimes \mu) \equiv \mu;(\mu \otimes \mathrm{id})\}$
    \vfill
    \end{example}
    \vfill
    Two kinds of composition: $\fatsemi$ and $\otimes$
    \vfill
    Additional generators $\id_n$, $\sym_{m,n}$
    \vfill
    Must coherently interact with each other: $f \fatsemi \id_n = f$, $(f \otimes g) \fatsemi (h \otimes k) = (f \fatsemi h) \otimes (g \fatsemi k)$, $\ldots$
\end{frame}

\begin{frame}{Symmetric monoidal theories}
    Can be still presented algebraically by lifting $\fatsemi$ and $\otimes$ to the level of functional symbols
    \vfill
    \begin{example}[$((\mu \fatsemi \sym) \otimes (\eta \otimes \id)) \otimes (\mu \fatsemi (\id \otimes \id))$]
    \vfill
\[    
\adjustbox{scale=0.6}{\tikzfig{./figures/e-graph-comonoid-example}}
\]
    \end{example}
\end{frame}

\begin{frame}
\vfill
\begin{itemize}
\item Saturating the e-graph yields...
\end{itemize}
\centering
\small
\begin{tabular}{@{}lrr@{}}
\toprule
Rule & $(\to)$ & $(\leftarrow)$ \\
\midrule
Left identity \hfill $(\id_n \fatsemi f = f)$ & --- & 899 \\
Right identity \hfill $(f \fatsemi \id_n = f)$ & 4 & 308 \\
Composition assoc.\ \hfill $((f \fatsemi g) \fatsemi h = f \fatsemi (g \fatsemi h))$ & 1{,}325 & 580 \\
Bifunctoriality \hfill $((f_1 \otimes f_2) \fatsemi (g_1 \otimes g_2) = (f_1 \fatsemi g_1) \otimes (f_2 \fatsemi g_2))$ & 157 & 92 \\
Tensor assoc.\ \hfill $((f \otimes g) \otimes h = f \otimes (g \otimes h))$ & 1{,}695 & 2{,}174 \\
Tensor unit \hfill $(\id_0 \otimes f = f,\; f \otimes \id_0 = f)$ & --- & 2{,}005 \\
Sym.\ naturality \hfill $(\sym \fatsemi (f \otimes g) = (g \otimes f) \fatsemi \sym)$ & 10 & --- \\
Sym.\ involution \hfill $(\sym \fatsemi \sym = \id)$ & --- & 1 \\
Commutativity \hfill $(\mu \fatsemi \sym = \mu)$ & 1 & --- \\
Counit \hfill $(\mu;(\eta \otimes \id) = \id)$ & 1 & --- \\
\midrule
Total & 3{,}193 & 6{,}059 \\
\bottomrule
\end{tabular}
\end{frame}

\begin{frame}{Variables and Binding}
Traditional e-graphs model variables promoting them to functional symbols with no inputs. No notion of binding and substitution
\vfill
\begin{example}[{[\cite{OptimizingBetaReductionEGraphs}]}]
\vfill
% \[
% \adjustbox{width=\textwidth}{
% \tikzfig{./figures/e-graph-substitution}
% }
% \]
\begin{minipage}{0.1\linewidth}
    \[
   \adjustbox{scale=0.55}{\tikzfig{../../../figures/categorical-semantics/e-graph-substitution-1}} 
    \]
\end{minipage}
\hspace{1em}
\begin{minipage}{0.2\linewidth}
    \[
   \adjustbox{scale=0.55}{\tikzfig{../../../figures/categorical-semantics/e-graph-substitution-2}} 
    \]
\end{minipage}
\hspace{2em}
\begin{minipage}{0.25\linewidth}
    \[
   \adjustbox{scale=0.55}{\tikzfig{../../../figures/categorical-semantics/e-graph-substitution-3}} 
    \]
\end{minipage}
\hspace{2em}
\begin{minipage}{0.25\linewidth}
    \[
   \adjustbox{scale=0.55}{\tikzfig{../../../figures/categorical-semantics/e-graph-substitution-4}} 
    \]
\end{minipage}
% \begin{minipage}{0.15\linewidth}
%     \[
%    \adjustbox{scale=0.6}{\tikzfig{../../../figures/categorical-semantics/e-graph-substitution-5}} 
%     \]
% \end{minipage}
% \hfill
\end{example}
\end{frame}

\section{Categorical E-graphs}

\begin{frame}{What we do}
Both challenges can be addressed by studying e-graphs from a \alert{categorical} perspective
\vfill
We view e-graphs as morphisms in a particular category induced by the signature
\vfill
Categorical approach has clear notions of variables and binding
\vfill
Helps to justify certain congruences stipulated by e-graphs
\vfill
A concrete representation of e-graphs can be obtained by considering \alert{string diagrams} for the corresponding category
\end{frame}
\begin{frame}{Categorical semantics~[\cite{tiurin2025equivalencehypergraphsdporewriting}]}
\vfill
To capture the notion of equivalence between (sub)terms we freely enrich a given signature-generated category over a category of semilattices
\vfill
Semilattices are sets with an idempotent commutative associative binary operation $+$
\vfill
This means we have the following properties in the category
\begin{itemize}
\item $f + f = f$
\item $f + g = g + f$
\item $(f + g) \fatsemi h = (f \fatsemi h) + (g \fatsemi h)$
\item $f \fatsemi (g + h) = (f \fatsemi g) + (f \fatsemi h)$
\item $f \otimes (g + h) = (f \otimes g) + (f \otimes h)$
\item $(f + g) \otimes h = (f \otimes h) + (g \otimes h)$
% \item $\Lambda_{A,B,C}(f + g) = \Lambda_{A,B,C}(f) + \Lambda_{A,B,C}(g)$
\end{itemize}
\vfill
\end{frame}

\begin{frame}{Or, string diagrammatically}
\[
\scalebox{0.8}{
\tikzfig{./figures/egraph-strings}
}
\]
\end{frame}


\begin{frame}{Or, string diagrammatically}
\[
\scalebox{0.55}{
\tikzfig{./figures/egraph-strings-equations}
}
\]
\end{frame}

\begin{frame}
    \begin{example}
        \begin{minipage}{0.4\linewidth}
\[    
\adjustbox{scale=0.8}{\tikzfig{./figures/e-graph-comonoid-example}}
\]
\end{minipage}
\hfill
\begin{minipage}{0.4\linewidth}
\[    
\adjustbox{scale=0.8}{\tikzfig{./figures/string_diagram_commutative_comonoid}}
\]
        \end{minipage}
    \end{example}
\vfill
A concrete representation can literally absorb all the non-domain-specific equations (and some domain-specific ones)~[\cite{bonchi-lics}]
\vfill
Extend this representation to support string diagrams with dashed boxes~[\cite{tiurin2025equivalencehypergraphsdporewriting}]
\end{frame}
\begin{frame}
    Traditional e-graphs are recovered as string diagrams for cartesian semilattice-enriched categories
    \begin{example}[$(a * 2) / 2$]
        \[
        \scalebox{0.8}{
        \tikzfig{./figures/e-graph-example-a-no-label}
        }\qquad\qquad \scalebox{0.8}{\tikzfig{./figures/e-string-example-a}}
        \]
    \end{example}
\end{frame}

\begin{frame}
\begin{example}
    \[
    \adjustbox{width=\linewidth}{
    \tikzfig{../../../figures/categorical-semantics/egraph-translation-step-by-step-a-b}
    }
    \]
\end{example}
\end{frame}

\begin{frame}

\begin{example}
    \begin{minipage}{0.25\linewidth}
\[
\adjustbox{width=\linewidth}{
\tikzfig{./figures/e-string-example-b}
}
\]
    \end{minipage}
\hfill
\begin{minipage}{0.05\linewidth}
    \[
    \xRightarrow{}
    \]
\end{minipage}
\hfill
\begin{minipage}{0.25\linewidth}
\[
\adjustbox{width=\linewidth}{
    \tikzfig{./figures/e-string-example-c}
}
\]
\end{minipage}
\hfill
\begin{minipage}{0.05\linewidth}
    \[
    \xRightarrow{}
    \]
\end{minipage}
\hfill
\begin{minipage}{0.25\linewidth}
    \[
    \adjustbox{width=\linewidth}{
    \tikzfig{./figures/e-string-example-d}
    }
    \]
    \end{minipage}
\end{example}
\end{frame}

\begin{frame}
\begin{example}
\begin{minipage}{0.35\linewidth}
        \[
        \adjustbox{width=\linewidth}{
        \tikzfig{./figures/e-string-example-d}
        }
        \]
\end{minipage}
\hfill
\begin{minipage}{0.2\linewidth}
    \[
    \xRightarrow{}
    \]
\end{minipage}
\hfill
\begin{minipage}{0.35\linewidth}
    \[
    \adjustbox{width=\linewidth}{
    \tikzfig{./figures/e-string-example-e}
    }
    \]
\end{minipage}
\end{example}
\end{frame}


\bsubsection{Combinatorial semantics}

\begin{frame}{This work}

Semilattice-enriched closed symmetric monoidal categories 
\vfill
\[
\adjustbox{width=0.9\textwidth}{
\tikzfig{./figures/egraph-strings-2}}
\]
\vfill
\begin{gather*}
f = \Lambda_{A,B,C}(f) \otimes \id_{B} \fatsemi \textbf{ev}_{B,C}\\
\id_{B \multimap{} C} = \Lambda_{B \multimap{} C,B,C}(\textbf{ev}_{B,C})\\
\Lambda_{Z,B,C}(g \otimes \id_{B} \fatsemi f) = g \fatsemi \Lambda_{A,B,C}(f)\\
\Lambda_{A,B,C}(f + g) = \Lambda_{A,B,C}(f) + \Lambda_{A,B,C}(g)
\end{gather*}
\end{frame}

\begin{frame}{Or, string diagrammatically}
\begin{minipage}{0.45\linewidth}
    \[
   \adjustbox{scale=0.6}{\tikzfig{./figures/beta_law}} 
    \]
\end{minipage}
\hfill
\begin{minipage}{0.45\linewidth}
    \[
   \adjustbox{scale=0.6}{\tikzfig{./figures/curry_eval}} 
    \]
\end{minipage}
\hfill
\vfill
\begin{minipage}{0.45\linewidth}
    \[
   \adjustbox{scale=0.6}{\tikzfig{./figures/g_id_f}} 
    \]
\end{minipage}
\hfill
\begin{minipage}{0.45\linewidth}
    \[
   \adjustbox{scale=0.6}{\tikzfig{./figures/curry_join}} 
    \]
\end{minipage}
\end{frame}

\begin{frame}
\begin{center}
$\lambda f . f\;((\lambda x . f \;x)\;2)$
\end{center}
\vfill
\[
\adjustbox{scale=0.6}{
\tikzfig{./figures/de-brujin-string}
}
\]
Invariant under $\alpha$-equivalence
\vfill
Invariant under non-interacting substitutions
\vfill
Substitution is composition
\vfill
$\beta$-reduction is diagram rewiring
\end{frame}

\begin{frame}
\[
\adjustbox{width=0.8\textwidth}{
\tikzfig{./figures/e-graph-substitution}
}
\]
\vfill
\[
\adjustbox{width=0.8\textwidth}{
\tikzfig{./figures/e-string-substitution-2}
}
\]
\end{frame}

\begin{frame}{Slotted e-graphs~[\cite{slotted-egraphs}]}
\[
\adjustbox{scale=0.6}{\tikzfig{./figures/slotted_1}} \qquad \Rightarrow{} \qquad \adjustbox{scale=0.6}{\tikzfig{./figures/slotted_2}}
\]
\end{frame}

\begin{frame}

Enrichment helps to justify certain equalities
\vfill
\begin{equation*}
	\begin{aligned}
		\Lambda_{A,B,C}(f + g) & = \eta_A;(B \multimap (f + g))                    \\
		                       & = \eta_A;(B \multimap f + B \multimap g)          \\
		                       & = \eta_A;(B \multimap f) + \eta_A;(B \multimap g) \\
		                       & = \Lambda_{A,B,C}(f\;) + \Lambda_{A,B,C}(g).
	\end{aligned}%
	\label{law:distributivity}
\end{equation*}
\vfill
That appeared in slotted e-graphs

\begin{equation*}%
    \label{eq:cong-bind}
    \inferrule*[right=cong. (bind)]
    {E \vdash tc \cong tc'}
    {E \vdash \text{bind}\; \$x\; tc \cong \text{bind}\; \$x\; tc'},
\end{equation*}
\end{frame}

\begin{frame}{Lambda calculus with explicit substitution aka let~[\cite{Standardisation}]}
\[t \coloneqq x \mid t t \mid \lambda x. t \mid t [x/t]\]
\vfill
\begin{align*}
	(\lambda x. t) L u              \quad & \to_{d\beta} \quad t[x/u] L,                                                \\
	C \llbracket x \rrbracket [x/u] \quad & \to_{ls} \quad C \llbracket u \rrbracket [x/u],                             \\
	t [x/u]                         \quad & \to_{gc} \quad t,                               & x \notin \textsf{fv} (t),
\end{align*}
\vfill
\begin{align}
	t [x/u] [y/v] \quad        & \sim \quad t [y/v] [x/u],        & x \notin \textsf{fv} (v) \land y \notin \textsf{fv} (u), \nonumber \\
	\label{eq:subst-lambda}
	(\lambda y. t) [x/u] \quad & \sim \quad \lambda y. (t [x/u]), & y \notin \textsf{fv} (u), \tag{$\ast$}                             \\
	(t v) [x/u] \quad          & \sim \quad t [x/u] v,            & x \notin \textsf{fv} (v), \nonumber
\end{align}
\end{frame}

\begin{frame}
    \adjustbox{width=\textwidth}{
	\begin{tabular}{lcccc}
		Set of equations & $d\beta$ + $ls$ + $gc$ & $d\beta$ + $ls$ + $gc$    & $d\beta$ + $ls$ + $gc$     & $d\beta$ + $ls$ + $gc$\\
		                 &                        & + \texttt{lambda\_subst}  & + \texttt{lambda\_subst}   & + \texttt{lambda\_subst}\\
		                 &                        &                           & + \texttt{subst\_comm}     & + \texttt{subst\_comm}\\
						 &                        &                           &                            & + \texttt{subst\_app}\\						 
		\# of e-nodes    & 88                     & 185                       & 244                        & 846
	\end{tabular}}
\end{frame}	
\begin{frame}
    \adjustbox{width=\textwidth}{ 
    \begin{tabular}{l|ccccc}
		Rewrite rule & \# of applications     & \# of applications         & \# of applications      & \# of applications       & Absorbed?\\
		             &($d\beta$ + $ls$ + $gc$)&($d\beta$ + $ls$ + $gc$     &($d\beta$ + $ls$ + $gc$  &($d\beta$ + $ls$ + $gc$   &\\
		             &                        & + \texttt{lambda\_subst})  & + \texttt{lambda\_subst}& + \texttt{lambda\_subst} &\\
		             &                        &                            & + \texttt{subst\_comm}) & + \texttt{subst\_comm}   &\\
		             &                        &                            &                         & + \texttt{subst\_app})   &\\
		\hline
		\texttt{beta} & 14 & 57 & 49 & 80 & No\\
		\texttt{subst\_no\_var} & 35 & 113 & 118 & 197 & No\\
		\rowcolor{gray!30}
		\texttt{ls\_rule\_*} & 48 & 141 & 140 & 302 & Yes\\
		\texttt{subst\_lambda} & --- & 76 & 109 & 336 & No\\
		\rowcolor{gray!30}
		\texttt{subst\_comm} & --- & --- & 104 & 321 & Yes\\
		\rowcolor{gray!30}
		\texttt{subst\_app} & --- & --- & --- & 263 & Yes\\
	\end{tabular}}
\end{frame}

\begin{frame}{Conclusion}
  We extend the categorical semantics of e-graphs to support \alert{binding},
  moving from semilattice-enriched symmetric monoidal to semillattice-enriched symmetric monoidal \alert{closed} categories
  \vfill
  A single string diagram encodes both \alert{binding} (closed structure) and
  \alert{equivalence} of subterms (semilattice enrichment)
  \vfill
  \alert{Closed e-hypergraphs} give these diagrams a combinatorial
  representation, with DPOI rewriting shown \alert{sound and complete} w.r.t.\
  term rewriting modulo SMC equations
  \vfill
  Binding comes \alert{for free}: no explicit-substitution nodes or de Bruijn
  shifting --- $\beta$-reduction is just re-wiring, and $\alpha$-equivalence built in
  \vfill These \emph{categorical e-graphs} have clear advantages over slotted e-graphs in some scenarios and built from a semantics-first approach
\end{frame}


\appendix

\begin{frame}[allowframebreaks]{References}
\footnotesize
\printbibliography
\end{frame}

\end{document}